\section{Hardware and Software} \label{hardware_and_software}
    Hardware is the electronic and electromechanical components of a computer, for example the CPU, the RAM or a monitor
    Software is code that runs on the hardware, for example apps or the operating system.
    Software does not have a phyisical presence.

    \subsection{Application Software}

    \subsection{System Software}
        System software is software that runs to support application software.

        \subsubsection{Operating Systems}
            An operating system is a middle-man between the user/software and the hardware.
            Its tasks include:

            \paragraph{Processor Scheduling}
            Processor scheduling is how the computer decides what processes should be done and in what order.
            It does this by assigning priorities to each of the processes, and then uses an algorithm to decide which to do first.

            The first-come first-serve algorithm (FIFO) immediately executes the first process that arrives, and all other processes join a first-in first-out queue structure to be executed in order.
            This can mean that high-priority processes or processes with short deadlines can be ignored for a long time if there are a lot of other processes already in the queue.

            The earliest deadline first algorithm (EDF) works by finding the processes with the soonest to expire deadlines and executes them first.

            The round robin algorithm allocates a unit of time to each process, and then cycles to the next one whether or not the first has finished, and then it will be cycled back to later.
            These units of time are called time slices.
            A time slice, AKA a quantum, is a short unit of time where processes can execute, allowing for multitasking.

            Fixed-priority preemptive scheduling is the algorithm used by Windows 11, and involves ranking each process based on priority.
            Factors such as deadlines can affect the priority of a process.
            The processes are then executed in order of priority, with any new high-priority processes jumping in front of the other, lower-priority ones.

            \paragraph{Memory Management}
            Memory management is a set of techniques used to keep a computer's memory organised and working as quickly as possible.
            This involves keeping track of what memory is in use and what memory is free.

            Paging uses a lookup table called a page table to map logical addresses (the number assigned to the memory location) to the physical location (the physical place the data is stored), meaning that it can be stored anywhere, and is not stored in continguous pieces.
            
            Segmentation stores data in continguous blocks based on their relations to each other, meaning that a program would be stored as one block of data in memory.
            A segment could store a program module, for example.

            Virtual memory is where parts of the secondary storage is used as RAM.
            This memory will be slower, since secondary storage is not as fast as RAM, however computers generally have more secondary storage than RAM, since it is much cheaper.
            As RAM prices go up, virtual memory is becoming more useful for cheaper computers that have less RAM.

            Disk thrashing is a performance drop caused by the computer swapping lots of data between secondary storage and RAM, which causes the CPU to slow down.
            This is in turn caused by bad memory management, which should aim to minimise disk thrashing by copying larger chunks of data from secondary storage to RAM, with the hope that the nearby data will also be useful.
            This can also be caused by not having enough RAM, which means that not enough data can be stored at once.

            Memory management means that the computer cannot overwrite memory that is in-use, which would cause programs to have to fetch data back from secondary storage.
            It also allows multiple programs to occupy the same memory, which can be allocated properly to avoid data being overwritten.
            It also means that programmers don't have to worry about where their programs and data are stored, since it is all handled by the OS.

            \paragraph{Backing Store Management}
                a

            \paragraph{IO Management}
                a

    \newpage
    
\section{Operating Systems} \label{operating_systems}

    \newpage

\section{Programming Language Classification} \label{programming_language_classification}
    \subsection{Machine Code}
        Machine code is the binary 1s and 0s that enter a computer's CPU and control what it does.
        Machine code is rarely written by humans, since it is very difficult to read and write, however the first programs, written as holes on punch cards were written directly in machine code.

    \subsection{Assembly}
        Assembly was created as a more readable form of machine code.
        Assembly uses opcodes and operands to make instructions.
        
        An opcode is a 3-letter mnemonic, short for a command, such as `store' (\verb|STR|).
        The opcode is then given operands, which are pieces of data that the command uses, for example, \verb|STR| takes a register \verb|Rn| and a memory reference, where it will store the contents of the \verb|n|th register.
        Operands are either a memory address (written as a decimal number), a register (written with a capital R before a decimal number), or a value (written with a hash before a decimal number).
        To store the value in register 2 into memory location 300, assembly would use the line:
        \begin{verbatim}
STR R2 300
        \end{verbatim}

        The full instruction set for AQA assembly language can be found \href{https://pmt.physicsandmathstutor.com/download/Computer-Science/A-level/Past-Papers/AQA/AS-Paper-2/Assembly\%20Language\%20Instruction\%20-\%20Paper\%202\%20AQA\%20Computer\%20Science\%20AS-level.pdf}{here}.
        This is an example of a program in assembly: % TODO: Check that this code actually works
        \begin{verbatim}
// Compute the 10th Fibonacci number
LDR R0 #0   // R0 will be the counter to check how many
            // iterations have happened
LDR R3 #0   // R1 is the first number to add; it will be
            // shifted into R3
LDR R1 #1   // R2 is the second number to add

fibonacci:          // A label to start the loop
    CMP R0 #9       // Check whether more than 9 iterations
                    // have happened yet
    BGT exit        // If yes, exit
    LDR R2 R1       // Shift R1 into R2
    LDR R1 R3       // Shift R3 (previous sum) into R1
    ADD R3 R1 R2    // Adds the two previous numbers
    ADD R0 R0 #1    // Increments the counter
    B fibonacci     // Loops back to the fibonacci label

exit:               // Exit label for when the program
                    // finishes
    STR R3 100      // Stores the result to memory location
                    // 100
    HALT            // Stop the program
        \end{verbatim}

        Assembly has a one-to-one correspondance with machine code.
        This means that it can be converted, one instruction mapping to one machine code instruction.
        For example, the command \verb|ADD R3 R1 R2| could be assembled into the machine code 0011 0011 0001 0010; the opcode (\verb|ADD|) assembling to 0011, and the registers assembling to 0011 (\verb|R3|), 0001 (\verb|R1|), and 0010 (\verb|R2|).
        This makes it very easy to convert into machine code, hence many early programs were written in it, but it makes it hard to read, even if it is easier than pure machine code.
        However, machine code gives the programmer complete control over the memory, down to each bit.
        This means that programs can be optimised perfectly for the CPU, and leave no wasted memory, whereas running code through higher-level languages can cause unnecessary machine code being added.

        Each processor has its own assembly instruction set, which will be converted to machine code that that processor understands.
        This means that assembly can only work on one type of processor, whereas high-level languages will work for any processor that has a suitable interpreter/compiler for it.

        Overall, these qualities make assembly good for highly specific programs that run on one chip, but have to run very quickly.
        For example, an integrated system might use programs written in assembly, since they will waste little to no memory, will only have to run on that one chip, and have to run efficiently.

    \subsection{High-Level Languages}
        A high-level language is a language that looks more like English than machine code or assembly.
        The first high-level languages were languages such as Fortran, created by John Backus to look more like English.
        The program above to compute the 10th Fibonacci number looks like this in Fortran:
        \begin{verbatim}
! Compute the 10th Fibonacci number
program fibonnaci
    integer :: n, i, a, b, sum  ! Declare all variables

    n = 10      ! Initisalise max iteration to 10
    a = 0       ! Set the first number to 0
    b = 1       ! Set the second number to 1

    do i = 2, n         ! Iterate from i = 2 to i = n
        sum = a + b     ! Add a to b
        a = b           ! Shift b to a
        b = sum         ! Shift sum to b
    end do

    print *, sum        ! Output the final sum
end program fibonnaci
        \end{verbatim}
        This shows how much easier to read high-level languages are over low-level languages such as assembly.
        Variables mean that values can have meaningful names rather than just numbers, and structures such as iteration (\verb|do i = 2, n|) allow programs to be much shorter by handling multiple assembly commands in one line.

        High-level languages mean that programs are easier to code, decreasing production time and making the code more maintainable for future programmers.

    \subsection{Interpreted Languages}
        Some high-level languages are interpreted, meaning that they are read by an interpreter program line by line, and executed line by line as they are read.
        This means that programs have to be reread every time they are executed, which will slow them down.
        An example of an interpreted language is Python.

        Interpreted languages require reinterpreting every time they are run, meaning that they can only be run if the person running it has an the interpreter and the source code, which means that the code can be easily copied by anyone with access to it.
        This makes it not ideal for large programs that are being sold.

    \subsection{Compiled Languages}
        Compiled languages are read in their entirety by a compiler, which compiles the source code (the code written by the developer) into a file of obeject code (machine code) called an executable.
        This can be directly run by a CPU, meaning that it is quick to run.
        However, any small edit to the code requires a full recompile, meaning that development can be slow for large projects that need a lot of recompiles.
        Some examples of compiled languages are VB.Net, C++, Java and Fortran.

        Compiled code can be shared easily as an executable, meaning that the source code doesn't have to be shared as well as the program, making it ideal for a lot of application software.

    \newpage